\scanpage{25}
% Pencil foliation: 24, lower left.
$M_{i_0}\in\operatorname{Ob}D^b(\mathcal{M}^+(X_{i_0}))$, d'où $M_i,M$, on
aura pour $f_{i_0}:X_{i_0}\unsure{-|\longrightarrow}Y_{i_0}$ la relation
\[
 Rf_*(M)=\varinjlim_i v_i^*(Rf_{i*}(M_i))
\]
où $v_i:Y\unsure{-|\longrightarrow}Y_i$ est le morphisme
canonique.

4) Soit $\ell\in\mathbb{P}$ tel que $\ell\,1_X\in\Gamma(X,\mathcal O_X)$ soit
inversible. Alors on a un foncteur
\diagram{\begin{tikzcd}
T_\ell=T_\ell^{(X)}:\mathcal{M}^+(X)\arrow[r]&\mathcal{M}_\ell^+(X)
\end{tikzcd}}
où $\unsure{\mathcal{M}_\ell^*(X)|\mathcal{M}_\ell^+(X)}$ est \inserted{la catégorie}
formée des ``$\mathbb{Q}_\ell$-modules
constructibles sur $X$'', i.e. la catégorie
déduite de la catégorie des ``systèmes
$\ell$-adiques de faisceaux de $\ell$-torsion
constructibles'' en négligeant précisément
les faisceaux de torsion.
Le foncteur
$T_\ell$ est \inserted{\uuline{compatible avec} $\otimes$ \uuline{et unité}}
\uuline{exact} et \fbox{\uuline{fidèle}} (mais non
pleinement fidèle), \uuline{compatible avec le}
